\documentclass[10pt,a4paper]{article}
\usepackage[T1]{fontenc}
\usepackage[utf8]{inputenc}
\usepackage{lmodern}
\usepackage{amsmath,amssymb}
\usepackage[margin=21mm]{geometry}
\usepackage{microtype}
\usepackage[hidelinks]{hyperref}
\usepackage{enumitem}
\setlist[itemize]{leftmargin=*,nosep,topsep=3pt}
\setlength{\parindent}{0pt}
\setlength{\parskip}{5pt}
\setlength{\emergencystretch}{3em}
\newcommand{\question}[1]{\par\medskip{\normalsize\bfseries #1}\par\smallskip}
\pagestyle{plain}

\begin{document}
{\Large\bfseries Reevaluating the role of the STN-GPe network in generating beta synchronized oscillatory activity in parkinsonism\par}

\textbf{Paper:} De la Crompe, Deffains, Leblois and Mallet (2025).\quad Handbook of Parkinson's Disease Mechanisms, chapter 17.\par
\href{https://doi.org/10.1016/B978-0-443-21992-4.00013-0}{DOI: \nolinkurl{10.1016/B978-0-443-21992-4.00013-0}}\par
{\footnotesize Page references below use PDF page numbers, including front matter.}

\question{To which general question does this paper relate?}
Origins of pathological synchronization and which circuit components should be modelled. Critically reassesses a homogeneous STN--GPe pacemaker hypothesis.

\question{How does it answer the general questions?}
\textbf{Evidence.} Synthesizes recordings and selective optogenetic perturbations, including distinctions between prototypic and arkypallidal GPe neurons. In the reviewed parkinsonian-rat experiments, GPe inhibition disrupted beta synchronization whereas STN or motor-cortex inhibition did not (PDF pp.~4--8).

\textbf{Usable insight.} Pallidostriatal interactions and population-specific connectivity can explain patterns missed by a two-population loop. Species and preparation matter. Motor deficits can precede exaggerated beta, so the chapter does not equate the origin of synchronization with the sole cause of neurological symptoms (PDF pp.~1, 8--10).

\textbf{Model relevance.} Motivates separating GPe populations and including striatal feedback when the experimental setting requires them. It supplies structural constraints and competing mechanisms, rather than one numerical notebook parameterization. A defensible control objective needs behavioural evidence in addition to a change in beta power.

\question{What papers does it cite that could provide more information?}
\begin{itemize}
\item De la Crompe et al. (2020), \emph{The globus pallidus orchestrates abnormal network dynamics in a model of Parkinsonism}: causal perturbation evidence (cited in Figs.~17.4--17.5, PDF pp.~7--8).
\item Azizpour Lindi, Mallet and Leblois (2024), \emph{Synaptic Changes in Pallidostriatal Circuits\ldots}: constrained alternative circuit model (ref.~48, PDF p.~12; note~20).
\end{itemize}

\end{document}
